\documentclass[10pt]{article}

\usepackage{parskip}
\usepackage[margin=1in]{geometry} 
\usepackage{amsmath,amsthm,amssymb, graphicx, multicol, array}
\usepackage{makecell}
\usepackage{enumitem}
\usepackage{amssymb}
\usepackage{float}
\usepackage{href-ul}
 
\newcommand{\N}{\mathbb{N}}
\newcommand{\Z}{\mathbb{Z}}
 
\newenvironment{problem}[2][Problem]{\begin{trivlist}
\item[\hskip \labelsep {\bfseries #1}\hskip \labelsep {\bfseries #2.}]}{\end{trivlist}}

\begin{document}
 
\title{Exam}
\author{Jakob Sverre Alexandersen\\
GRA4154 Supply Chain Optimization with Mathematical Programming\\
2025}
\maketitle

\tableofcontents
\newpage
\section{Q 1}
Model used to optimize a manufacturing company's supply network structure before entering a new market (new country). The index $i $ represents candidate locations for factories in the country. The index $j $ represents customers who will receive products from the company's factories. 

\textbf{Model}
\begin{gather}
    \min Z = \sum_i\sum_j v_{i, j}X_{i, j} + \sum_i f_iY_i\notag\\
    \text{s.t.}\notag\\
    \sum_j X_{i, j} \le \text{cap}_i \times Y_i\quad\forall i\\
    \sum_i X_{i, j} = d_j\quad\forall j\\ 
    X_{i, j} \le d_j\times Z_{i, j} \quad\text{for each $i $ combination of $i $ and $j $}\\
    J\times W_j \ge \sum_i Z_{i, j} - 1 \quad\forall j\\ 
    \sum_j W_j \le 2\\
    X_{i, j} \ge 0 \quad\text{for each $i $ combination of $i $ and $j $}\\
    Y_i = 0 \text{ or } 1 \quad\forall i\\
    Z_{i, j} = 0 \text{ or } 1 \quad\text{for each $i $ combination of $i $ and $j $}\\
    W_j = 0 \text{ or } 1 \quad\forall j
\end{gather}
\begin{enumerate}
    \item Give a brief explanation of the objective function: 

    We minimize where $v_{i, j} $ is the variable cost of shipping products from factory $i $ to customer $ j $, times the amount of items to be shipped $X_{i, j} $. The other sum is the fixed cost $f_i $ of setting up the factory at location $i $ times the binary variable $Y_i $ (whether they want to set up a factory there or not). 
    \item What does constraint (2) achieve?

    This one sets supply to the same as demand. 
    \item Give a brief explanation of the variable $Z_{i, j} $

    It's a (binary) decision variable indicating whether any product is shipped from $i $ to $j $
    \item Why is constraint (3) in the model?

    It connects the binary variable $Z_{i, j} $ to the continuous variable $X_{i, j} $ such that no quantity is shipped from factory $i $ to market $j $ unless the binary variable is equal to 1
    \item Why is constraint (4) in the model?

    $W_j $ may be an binary variable signifying whether a customer has more than one source. Thus, we force $W_j $ to be 1 if there are more than one $X_{i, j} $ for that market. $J $ is the big M here
    \item What does constraint (5) achieve?

    It ensures that at most 2 customers get deliveries from multiple sources. 
    \newpage
    \item The above model is an example of a MILP model where all multiplications consist of a parameter being multiplied with a variable. However, the same logical relations could also be formulated as a MI\underline{N}LP model where some of the variables are multiplied with some of the other variables. What is the advangate of choosing a MILP model formulation instead of a MINLP model formulation?

    Both MILP and MINLP are NP-hard, but MILP gives accessibility to strong relaxations, and therefore good pruning, and they're primarily solved with branch-and-bound. The problems are fast, convex and globally solvable. MINLP, on the other hand, are either convex, or non-convex. Convex MINLP problems are still harder than MILP because each node may require solving an NLP, whereas the non-convex case is a different animal. It requires global optimization, which results in a combinatorial explosion. 
    \item Add the following logic: 

    Iff. products are shipped from a given factory (factory $D $) to a given customer ($7 $), then the quantity shipped from $D $ to 7 must be above 3,000 tons. 

    \begin{gather*}
        X_{D, 7} \ge 3000 \times Z_{D, 7} 
    \end{gather*}
    \item Add the following logic: 

    Shipping from factory $B $ to customer 12 would require some technical adjustments at factory $B $. Those adjustments would imply that the total capacity of factory $B $ is reduced by 600 tons 

    \begin{gather*}
        \sum_j X_{B, j}\le \text{cap}_B\times Y_B - 600\times Z_{B, 12}
    \end{gather*}
    \item Add the following logic: 

    If both factory $E $ and $G $ are built, then factory $H $ cannot be built
    \begin{gather*}
        Y_H \le 2 - Y_E - Y_G
    \end{gather*}
\end{enumerate}
\newpage 
\section{Q 2}
A company needs to select a portfolio of investment projects to maximize the total NPV, given an investment budget constraint. All projects are of the type ``all in or nothing'', i.e., it is not possible to invest in a fraction of a project. 

Show how the optimal solution is found by solving each step of the branch and bound method manually. What is the optimal portfolio of investment projects?

\begin{center}
    \begin{tabular}{|c|c|c|c|c|}
        \hline 
        \textbf{Project} & \textbf{A} & \textbf{B} & \textbf{C} & \textbf{D}\\\hline 
        \textbf{Value} & 35 & 40 & 21 & 64 \\\hline 
        \textbf{Investment} & 20 & 30 & 15 & 40 \\\hline 
        \textbf{Value per investment} & 1.75 & 1.33 & 1.4 & 1.6 \\\hline 
    \end{tabular}
\end{center}
\begin{gather*}
    \max Z = 35a + 40b + 21c + 64d\\
    \text{s.t.}\\
    20a + 30b + 15c + 40d \le 50 \\
    a, b, c, d \in \{0, 1\}
\end{gather*}
Optimal solution to LP relaxation: $a = 1, b = 0, c = 0, d = 0.75 $. Branch on $ d $:
\begin{enumerate}
    \item $d = 1 $

    $a = 0.5, d = 1 $: UB = 81.5: infeasible, branch on $a $
    \begin{enumerate}
        \item $a = 1 $

        $a = 1, d = 1 $: infeasible, no branch 
        \item $a = 0 $

        $c = 0.66, d = 1 $: UB = 78, infeasible, branch on $c $
        \begin{enumerate}
            \item $c = 1 $

            $d = 1, c = 1 $: infeasible, no branch 
            \item $c = 0 $

            $d = 1, a = 0, c = 0, b = 0.33 $: UB = 77.2, infeasible, branch on $b $
            \begin{enumerate}
                \item $b = 1 $

                $d = 1, a = 0, c = 0, b = 1 $: infeasible, no branch 
                \item $b = 0 $

                $d = 1 $: UB = 64, no branch 
            \end{enumerate}
        \end{enumerate}
    \end{enumerate}
    \item $d = 0 $

    $a = 1, b = 0.5, c = 1 $: UB = 76: infeasible, branch on $b $
    \begin{enumerate}
        \item $b = 1 $

        $a = 1, b = 1 $: UB = LB = 75, feasible, no branch. \textbf{Optimal solution}
        \item $b = 0 $

        $a = 1, c = 1 $: UB = 56, feasible, no branch 
    \end{enumerate}
\end{enumerate}

\newpage
\section{Q 3}
A Norwegian rental company rents out heavy equipment and machinery to construction companies.

The customers rent the equipment for a given number of weeks. At the end of the rental
period, the equipment is returned to the rental company and will need some repair work
before it can be rented to other customers or kept in inventory waiting for demand to occur.

The rental company can cover demand either from its inventory of repaired equipment, or
by purchasing new equipment from external suppliers.

Due to limited repair capacity, there will often be an inventory of returned equipment that is
not yet repaired. For each unit of equipment that needs to be stored from one week to the
next, there is an inventory holding cost per unit per week.

For both unrepaired and new/repaired equipment, the storage capacity is limited. For
unrepaired equipment, there is a joint storage capacity (i.e., the total capacity for all types of
unrepaired equipment). For new/repaired equipment, there is a separate capacity for each
type of equipment.

\textbf{Parameters}
\begin{itemize}
    \item $d_{j, t} $: Forecasted demand for equipment $j $ in week $t $
    \item $f_{j, t} $: Estimated quantity of equipment type $j $ returned from customers in week $t $
    \item $r_{t, j} $: Repair time (hrs) per unit of equipment type $j $ repaired in week $t $
    \item $\text{cap} $: Total number of hours available for repair in week $t $
    \item $pc_{j, t} $: Purchasing cost per unit of equipment $j $ purchased in week $t $
    \item $rc_{j, t} $: Repair cost per unit of equipment $j $ repaired in week $t $
    \item $uhc_j $: Inventory holding cost per unit per week for for \underline{un}repaired equipment on inventory
    \item $nhc_j $: Inventory holding cost per unit per week for \textbf{new} or \textbf{repaired} equipment on inventory
    \item $a $: Joint storage capacity for unrepaired equipment
    \item $b_j $: Storage capacity for equipment type $ j $ (new or repaired)
\end{itemize}
For each week within a given planning horizon, the company wants to determine the optimal
number of units to purchase and the optimal number of units to repair. The goal is to
minimize the sum of purchasing costs, repair costs, and inventory holding costs.

\textbf{Variables}
\begin{itemize}
    \item $X_{j, t} $: Number of units to purchase of type $j $ in the beginning of week $ t $
    \item $Z_{j, t} $: Nuber of units to repair of type $j $ in week $ t $
    \item $H_{j, t} $: Number of new/repaired units of type $j $ on inventory in the end of week $t $
    \item $S_{j, t} $: Number of unrepaired units of type $j $ on inventory in the end of week $t $
\end{itemize}

Build a model to minimize the sum of purchasing costs, repair costs, and inventory holding costs.
\newpage 
\textbf{Model}
\begin{gather}
    \min Z = \sum_{j}\sum_tX_{j, t}\times pc_{j, t} + Z_{j, t}\times rc_{j, t} + S_{j, t}\times uhc_j + H_{j, t}\times nhc_j\notag\\
    \text{s.t.}\notag\\
    H_{j, t} = H_{j, t - 1} + X_{j, t} + Z_{j, t} - d_{j, t} \quad\forall j, t\\
    S_{j, t} = S_{j, t - 1} - Z_{j, t} + f_{j, t}\quad\forall j, t\\
    \sum_j rt_j\times Z_{j, t} \le \text{cap}\quad\forall t\\
    \sum_j S_{j, t} \le a\quad\forall t\\
    H_{j, t} \le b_j\quad\forall j, t\\
    X_{j, t} \ge 0\quad\forall j, t\\
    Z_{j, t} \ge 0\quad\forall j, t\\
    H_{j, t} \ge 0\quad\forall j, t\\
    S_{j, t} \ge 0\quad\forall j, t
\end{gather}
\begin{enumerate}
    \setcounter{enumi}{9}
    \item Ending inventory this week = inventory from last week + new units purchased + repaired units completed minus demand served 
    \item Ending inventory of unrepaired units this week = broken units from last week minus the amount of units to repair plus the estimated quantity of units to be returned 
    \item The total hours spent on repair time cannot exceed the shop's repair capacity
    \item The total amount of unrepaired units cannot exceed the capacity of the unrepaired units' storage 
    \item For each type of unit, the amount of units cannot exceed the capacity of that type's storage 
    \item Non-negativity constraints 
\end{enumerate}

\end{document}